% P1 学术写作版（2026-11-24）：结论与展望。由 cellcharter_study_conclusion_module_rev_p0_20261124.tex 改写；
% 局限的完整表述见第4.7.4节，此处只作概括。未编译。
\section{结论与展望}
\label{sec:cc-conclusion-p1}

\subsection{主要结论}
本研究以CellCharter原作者提出的方法和基准配置为基础，在3名供体的12张DLPFC切片和6例乳腺癌样本上，从邻域表征、聚类数规则、评价指标和跨组织适用性四个方面，系统评价了多尺度空间聚类的性能。主要结论如下。

\begin{enumerate}
  \item \textbf{邻域收益集中于最近邻层。}引入第一层邻域后，ARI由0.305升至0.408，并在全部切片中方向一致；15维之后的增益小于切片间差异，约半数切片和一名供体方向相反。不同15维邻域组合之间没有稳定占优者；在开发切片上提出的远邻替换组合，在其余切片和外部队列中没有复现其优势。在原作者的基准配置下，本研究结果与原作者发布的单切片结果处于同一水平。
  \item \textbf{聚类数规则对结论的影响大于深层尺度差异。}固定K时尺度排序稳定，但K本身引起的ARI变化远大于15维之后的尺度差异。参考类别数并非ARI最优的K。逐切片的无标签稳定性选K（5次重复）选出的簇数离散，并改变了尺度排序。
  \item \textbf{整体一致性与严格边界定位存在评价权衡。}随邻域加深，2-hop容差边界F1由0.286升至0.533，严格边界F1则从约0.145降至0.118。这一差异主要来自预测边界从碎片化走向平滑时召回率的变化，说明边界评价需要同时报告严格指标、容差指标及其分解。
  \item \textbf{方法比较与跨组织适用性依赖于条件。}在参考类别数条件下，CellCharter优于非空间方法和BANKSY，与BayesSpace相当；在K-blind条件下，比较结论受到各方法选K规则（包括本研究为BayesSpace和GMM设定的似然规则频繁命中上界）的强烈影响。在乳腺癌队列中，空间邻域特征相对非空间表示的收益在全部患者中方向一致，但患者间差异很大，与BayesSpace的比较随指标和K规则而变。
\end{enumerate}

综上，本研究表明：对于成熟的空间聚类方法，性能判断不能只依赖参考类别数条件下的单一整体指标，而应同时考察聚类数敏感性、严格与容差边界指标，以及供体或患者层面的结果。这为空间聚类方法的参数选择、结果解释和评价报告提供了可复核的实证依据。

\subsection{局限}
本研究的独立生物学样本较少（3名供体、6名患者），以描述性分层统计为主。部分配置所用的切片在调参或开发阶段已被分析过。邻域深度与特征维数同时变化。稳定性选K的重复次数少于原作者设定，部分基线的K-blind规则存在缺陷，对照集也未包含STAGATE。若干解释性分析属于事后分析。上述限制的完整说明见第\ref{subsubsec:cc-limits-p1}节。

\subsection{展望}
后续工作可以从以下几方面展开：
\begin{itemize}
  \item 在更多独立供体、组织类型和测序平台上检验当前结论的适用范围；
  \item 按各方法官方建议的选K方式（例如BayesSpace的拐点判断）及足够的稳定性重复次数，重新设计无标签条件下的对照，并纳入STAGATE等更强的基线；
  \item 以维度匹配的对照进一步区分邻域范围与特征维数的作用；
  \item 为边界指标建立零模型基准，并结合层间过渡区的标注不确定性完善边界评价。
\end{itemize}
